%% kapa: exam=1 round=26 year=2015 period=1 form=A subject=economics subject_ko=경제원론 seq=9 number=49
%% source: Q-Net 감정평가사 기출문제 (cst00309) artlSeq=5138001 file=extracted/1008/5138001/제26회_감정평가사 1차_기출/제26회_감정평가사 1차_기출/1교시_감평_A형.pdf
%% conversion: kordoc 4.24.51 -> tools/gen1.py (auto); review=proofread-sonnet; flags=none
\begin{question}{49}
甲기업의 단기 총비용함수가 $C=100+10Q$일 때, 甲기업의 비용에 관한 설명으로 옳지 않은 것은? (단, $Q$는 양(+)의 생산량이다.)
\begin{choices}
  \item 고정비용은 100이다.
  \item 모든 생산량 수준에서 한계비용은 10이다.
  \item 생산량이 증가함에 따라 총비용은 증가한다.
  \item 생산량이 증가함에 따라 평균비용은 감소한다.
  \item 모든 생산량 수준에서 한계비용은 평균비용보다 크다.
\end{choices}
\end{question}
\answer{5}
